% % Pages
% \usepackage{fullpage}

% % Margins
% \usepackage[top=2.5cm, left=3cm, right=3cm, bottom=4.0cm]{geometry}

% For math and formulas
\usepackage{amsmath,amssymb, newtxtext, newtxmath}

% Text and Font encoding
\usepackage[T1]{fontenc}
\usepackage[english]{babel}
\usepackage{bm}
\usepackage{csquotes}

% font format
\usepackage[scaled=0.86]{helvet}
\renewcommand{\familydefault}{\sfdefault}

% \usepackage{crimson}
% \renewcommand{\familydefault}{\sfdefault}

% Set spacing
\usepackage{setspace}
\setstretch{1}

% Numbering
\usepackage{enumitem}

% % paragraph styling
% \usepackage{parskip}

% % set paragraph indentation
\setlength\parindent{12pt}

% Images
\usepackage{graphicx}
% \usepackage{caption}

% Tables
% \usepackage[xcdraw]{xcolor}
\usepackage[table,x11names]{xcolor}
\usepackage{booktabs}

% % text color
% % \color{black}

% other 
\usepackage{comment} 
\usepackage{lipsum}

% \makeatletter
% % subsection numbering
% \renewcommand\thesubsection{
% \arabic{subsection})
% }

% % we use \prefix@<level> only if it is defined
% \renewcommand{\@seccntformat}[1]{%
%   \ifcsname prefix@#1\endcsname
%     \csname prefix@#1\endcsname
%   \else
%     \csname the#1\endcsname\quad
%   \fi}
% % define \prefix@section
% \newcommand\prefix@section{Problem \thesection: }
% \makeatother

% You might want to define your own abbreviated commands for commonly used terms, e.g.: 
\newcommand{\kms}{km\,s$^{-1}$}
\newcommand{\msun}{$M_\odot}

% % Section sizes
% \usepackage[small]{titlesec}

% Add the date in the title
\usepackage{etoolbox}
\patchcmd{\MaketitleBox}{\footnotesize\itshape\elsaddress\par\vskip36pt}{\footnotesize\itshape\elsaddress\par\parbox[b][36pt]{\linewidth}{\vfill\hfill\textnormal{\today}\hfill\null\vfill}}{}{}%
\patchcmd{\pprintMaketitle}{\footnotesize\itshape\elsaddress\par\vskip36pt}{\footnotesize\itshape\elsaddress\par\parbox[b][36pt]{\linewidth}{\vfill\hfill\textnormal{\today}\hfill\null\vfill}}{}{}%
